

% ============================================
% TÓM TẮT TIẾNG ANH (ABSTRACT - ENGLISH)
% ============================================

\newpage
\thispagestyle{empty}
\begin{center}
    \textbf{\Large ABSTRACT}
\end{center}

\vspace{0.5cm}

\noindent\textbf{Project Title:} Peer-to-Peer Electric Motorbike Sharing Platform for Local Communities

\noindent\textbf{Students:} Duong Hoang Long, Hang Nhut Long, Au Nguyen Hung Manh

\noindent\textbf{Supervisor:} PhD. Trương Tuấn Anh

\noindent\textbf{Co-supervisor:} Dương Huỳnh Anh Đức

\noindent\textbf{Academic Year:} 2024 - 2025

\vspace{0.5cm}

Vietnam currently has the highest motorcycle density in Southeast Asia, with over 85\% of personal transportation being two-wheeled vehicles. As the government promotes the transition to electric motorcycles to reduce air pollution and carbon emissions, the cost of electric vehicle ownership remains a significant barrier for many users, particularly students and urban residents. Meanwhile, personal vehicles in communities typically remain idle 70-80\% of the time, leading to substantial resource waste.

This project researches and develops a peer-to-peer (P2P) electric motorcycle sharing platform for local communities such as residential complexes, dormitories, and university campuses. Unlike centralized (B2C) models such as Grab or VinFast, the platform enables residents to directly share their personal vehicles with each other, optimizing idle asset utilization and reducing transportation costs.

The system is built on Clean Architecture combined with modern design patterns (Repository, Strategy, Observer, Factory), ensuring high modularity, maintainability, and scalability. 
The backend utilizes NestJS (TypeScript) with PostgreSQL and Prisma ORM for spatial data processing, Redis for caching and real-time location tracking. The frontend is developed using Flutter with BLoC pattern for state management, supporting cross-platform iOS and Android deployment.

Key system features include: 

(1) GPS-based vehicle search with multi-criteria filtering, 

(2) Booking and payment flow integrated with VNPay/MoMo payment gateways, 

(3) Real-time vehicle location updates via WebSocket, 

(4) Machine learning-based trust score system for user credibility assessment, 

(5) Two-way authentication and dispute resolution mechanisms to ensure transaction safety.

% The project has been piloted in 2-3 communities with 30-50 vehicles, 100-150 users, and 300-500 transactions over a 6-month period. Results demonstrate stable system operation with average response time under 2 seconds for primary operations, achieving 99.5\% uptime. Users provide positive feedback on convenience (4.2/5) and cost savings (40-50\% reduction compared to vehicle ownership).
\newpage
\thispagestyle{empty}
The main contributions of this project include: 

(1) Design and implementation of the first P2P sharing platform for electric motorcycles in Vietnam, 

(2) Application of Clean Architecture and design patterns to a real-world system, 

(3) Integration of machine learning for trust score calculation, 

(4) Resolution of challenges in real-time location tracking and concurrency control in booking systems, 

(5) Proposal of a business model adapted to Vietnamese culture and user habits.

The project contributes to promoting the sharing economy, reducing community transportation costs, and supporting the government's green transition objectives. Future development directions include expanding to multiple cities, integrating IoT for smart vehicles, and applying blockchain for transparent transaction history.

\vspace{0.5cm}

\noindent\textbf{Keywords:} Peer-to-Peer, Sharing Economy, Electric Motorcycle, Clean Architecture, Flutter, NestJS, Machine Learning, Trust Score, Real-time System

% ============================================
% END OF FRONT MATTER SECTIONS
% ============================================